\documentclass[11pt,a4paper]{article}

% --------------------
% PAQUETES
% --------------------
\usepackage[spanish]{babel}
\usepackage[utf8]{inputenc}
\usepackage[T1]{fontenc}
\usepackage{geometry}
\usepackage{enumitem}
\usepackage{hyperref}
\usepackage{xcolor}
\usepackage{titlesec}

\geometry{margin=2cm}
\setlength{\parindent}{0pt}

% --------------------
% ESTILO
% --------------------
\definecolor{sectioncolor}{RGB}{40,40,40}

\titleformat{\section}
  {\large\bfseries\color{sectioncolor}}
  {}
  {0pt}
  {}

\renewcommand{\labelitemi}{--}

% --------------------
% DOCUMENTO
% --------------------
\begin{document}

% --------------------
% ENCABEZADO
% --------------------
\begin{center}
    {\LARGE \textbf{JAVIER JOSÉ LOBATO CEBALLOS}}\\
    \vspace{2mm}
    \textbf{Ingeniero Electrónico}\\
    \vspace{2mm}
    \href{mailto:javierlobato@hotmail.com}{javierlobato@hotmail.com} \,|\, 
    +57 300 127 6123 \,|\, CC: 1004360307 \,|\, https://github.com/lob117
\end{center}

\vspace{4mm}
% --------------------
% PERFIL PROFESIONAL
% --------------------
\section*{Perfil Profesional}

Ingeniero Electrónico con enfoque en automatización de procesos e Innovación Digital, análisis de datos y tecnologías IoT.
Experiencia en el desarrollo de soluciones basadas en Inteligencia Artificial (OCR y geoparsing),
integración de flujos de trabajo automatizados mediante \textbf{n8n} y contenedores \textbf{Docker}.
Competente en el manejo de bases de datos \textbf{SQL} y lenguajes como \textbf{Python y C}.
Aplicados a la toma de decisiones estratégicas, captura de métricas operativas y optimización de procesos bajo metodologías ágiles.

% --------------------
% HABILIDADES
% --------------------
\section*{Habilidades Técnicas}

\begin{itemize}[leftmargin=*]
    \item \textbf{Ecosistema de Automatización:} Diseño de flujos en n8n y Power Automate.
    \item \textbf{Arquitectura de Datos e IA:} Gestión de bases de datos SQL, implementación de modelos RAG (Ollama/LangChain) y despliegue en contenedores Docker.
    \item \textbf{Innovación y Prototipado:} Desarrollo de Pruebas de Concepto (PoC) rápidas , integración de hardware IoT y 
    \item \textbf{Análisis de Procesos y datos:} Identificación de oportunidades de mejora, levantamiento de requerimientos técnicos, metodologías metodologías ágiles Scrum y informes técnicos (Excel y Google sheets).
    \item \textbf{Lenguajes:} Python, C/C++ y SQL.
\end{itemize}

% --------------------
% EXPERIENCIA
% --------------------
\section*{Experiencia Profesional}

\textbf{Corporación NaturalSIG} \\
\textit{Ingeniero Electrónico y Automatización de Procesos} \hfill Agosto 2025 -- Diciembre 2025

\begin{itemize}[leftmargin=*]
 \item Lideré la \textbf{ideación y diseño de soluciones} de automatización mediante el uso de OCR y Geoparsing, creando una gobernanza  de datos digitales.
    \item Diseñé la arquitectura funcional de flujos de trabajo en \textbf{n8n y Docker}, integrando componentes de IA para la transformación de datos no estructurados en información estratégica.
    \item Desarrollé prototipos rápidos (\textbf{PoC}) para validar la viabilidad técnica de integraciones entre plataformas GIS y bases de datos SQL.

\end{itemize}

\vspace{2mm}

\textbf{Corporación NaturalSIG} \\
\textit{Especialista en Datos UAV y Flujos Geoespaciales} \hfill Enero 2025 -- Junio 2025

\begin{itemize}[leftmargin=*]
    \item Gestión técnica y mantenimiento preventivo de plataformas UAV (EbeeX y Phantom 4).
    \item Automaticé la captura y procesamiento de datos masivos provenientes de sensores multiespectrales, optimizando el procesamiento de entrega de modelos 3D y mosaicos RGB.
    \item Aseguré la \textbf{integridad de las métricas} de rendimiento de los equipos UAV, gestionando el mantenimiento preventivo basado en análisis de datos operativos.
\end{itemize}


\vspace{2mm}

\textbf{Corporación NaturalSIG} \\
\textit{Practicante de Ingeniería (Optimización y Visión Artificial)} \hfill Agosto 2024 -- Diciembre 2024

\begin{itemize}[leftmargin=*]
    \item Desarrollé algoritmos de visión artificial para la automatización del conteo vehicular, transformando capturas de video en dashboards de análisis de tráfico.
    \item Diseñé estructuras de líneas eléctricas y edificaciones en AutoCAD.
\end{itemize}

% --------------------
% EDUCACIÓN
% --------------------
\section*{Educación y Certificaciones}

\textbf{Ingeniería Electrónica} \\
Universidad del Magdalena \hfill 2020 -- 2025

\textbf{Piloto de Dron Certificado (RAC 100):} APD (Idoneidad: 2449033). \\
APD \hfill 2025

\textbf{Inglés Conversacional (B1)} \\
IFT Cajamag \hfill 2025 -- 2026


% --------------------
% INFORMACIÓN ADICIONAL
% --------------------
\section*{Información Adicional}

\textbf{Idiomas:} Inglés B1 (Conversacional). \\
\textbf{Tarjeta Profesional:} AT206-188296. \\
\textbf{Herramientas:} Git/GitHub, VS Code, WSL2, Entornos Linux.

% --------------------
% REFERENCIAS
% --------------------
\section*{Referencias}

\textbf{Jaime Rodríguez Curcio} \\
Ingeniero Ambiental \\
\href{mailto:ing.ambiental1@naturalsig.org}{ing.ambiental1@naturalsig.org} \\
Cel: 302 408 6564

\vspace{2mm}

\textbf{Gustavo Adolfo Manjarres Pinzón} \\
Biólogo \\
\href{mailto:director@naturalsig.org}{director@naturalsig.org} \\
Cel: 321 772 1052

\end{document}
